\part[Computational behavioral characterization for the ecological assessment of dysexecutive syndromes]{Computational Behavioral Characterization for the \\Ecological Assessment of Dysexecutive syndromes}
\label{part:part_2}
% \part[Data-driven estimation of egocentric vision ethograms for the ecological assessment of dysexecutive syndromes]{Data-driven Estimation of \\Egocentric Vision Ethograms for the \\Ecological Assessment of Dysexecutive syndromes}


\chapter[Assessment of dysexecutive syndromes, study description, data collection]{Assessment of Dysexecutive syndromes, Study Description, Data Collection}
% \chapter[Definitions, study description, data collection, and multimodal outputs extraction]{Study description, data collection, and multimodal outputs extraction}
\label{sec:chapter_4}

% \begin{chapterabstract}
%     \textbf{Key points:}
%     \begin{enumerate}
%         \item Evaluating dysexecutive syndromes (DS) in brain-damaged patients remains a significant challenge in neurology, with scarce knwoledge of the mapping between brain lesions and their behavioral expression. Gold standard neuropsychological tests, although internally valid, often lack ecological validity, failing to reflect the difficulties patients report in everyday life.
        
%         \item This study aims to address this limitation,  improve care for individuals with dysexecutive syndromes, and map brain lesions to cognitive and behavioral impairments in real-world contexts,  by developing a multi-instrumented cooking test while the patient's behavioral and oculometric responses are recorded via three fixed cameras and wearable eye-tracking glasses recording egocentric vision.

%         \item During this test, three cameras film the test room, while the participants - equipped with oculometric glasses recording egocentric video flow, eye movements and head movements - perform a structured functional task consisting in following the steps of a chocolate cake recipe. 

%         % \item The Smartflat study develop a multi-instrumented cooking test administered in clinics, while the patient's behavioral and oculometric responses are recorded via three fixed cameras and wearable eye-tracking glasses recording egocentric vision.  

%         \item This part aim to leverage the large-scale (N=168) multimodal data collected during the Smartflat study to develop a computational behvaioral phenotyping approach to assess sensorimotor and executive dysfunctions. The spontaneous, diverse, and semantically complex nature of DS-related behavioral alterations support the development of data-driven, agnostic approaches with minimal supervision
        

\begin{chapterabstract}
    \textbf{Context and Motivation:}
    \begin{enumerate}
        \item Dysexecutive syndromes (DS) in brain-damaged patients remains a significant challenge in neurology, due to the complex and heterogeneous behavioral manifestations of executive dysfunction, and the limited understanding of the mapping between brain lesions and their behavioral expression. Gold standard neuropsychological tests, although internally valid, often lack ecological validity, failing to capture the difficulties patients report in their everyday life. 

        \item To address this gap, the Smartflat study introduces an ecologically valid, multi-instrumented behavioral assessment. Under examiner supervision, participants perform two structured tasks — preparing a chocolate cake and constructing the "Arc de triomphe" with Lego blocks — while their behavior is continuously monitored using three fixed cameras and wearable eye-tracking glasses that capture egocentric video, eye movements, gaze patterns, and head movements.

        \item This chapter initiates the development of a computational behavioral phenotyping pipeline for the analysis of the long-standing and large-scale Smartflat dataset. The spontaneous, diverse, and semantically complex nature of DS-related behavioral alterations support the development of agnostic, data-driven computational approaches with minimal supervision.

        % to improve the understanding of executive dysfunction by collecting multimodal data from ecological and semi-ecological tests. These tests are designed to be more reflective of real-world challenges faced by individuals with executive dysfunction.
        %\item The Smartflat study seeks to improve care for individuals with dysexecutive syndromes and to map brain lesions to cognitive and behavioral impairments in real-world contexts.

        %\item We aim in this part to leverage the large-scale (N=168) multimodal data collected during the Smartflat study to develop a digital phenotyping approach to assess sensorimotor and executive dysfunctions across the following five dimensions: static balance, dynamic balance, and locomotion; neurovegetative and emotional state; fine motor skills; visual information processing; and planning, initiation, mental flexibility, and memory. %These syndromes can result from various neurological conditions, including traumatic brain injury (TBI) and lesions in the prefrontal cortex (LPR).        

    \end{enumerate}

    \textbf{Technical Contributions:}
    \begin{enumerate}
        \item Development of a comprehensive multimodal data processing pipeline to analyze sensorimotor and executive dysfunctions, including data pre-processing, cleaning, management, and multimodal feature extraction from multi-view videos, audio, and eye-tracking. Two large-scale datasets are processed: the Cooking Task (CT) and the Arc de Triomphe Construction Task (ATCT), comprising 162 participants, 373 total administrations, with 1230 hours of multimodal data.

        \item Latent representations of short video segments ($\approx 640$ ms) using state-of-the-art video encoder allow the computation of symbolic representations of participants' egocentric vision trajectories, encoding behavior and task execution progress accurately. These aim to facilitate multimodal integration of hand and skeleton pose landmarks from multiple views, head movements, speech and speaker diarization, and fine-grained oculometric and pupillometric data, which altogether support a rich and comprehensive assessment of sensorimotor and executive dysfunctions from behavior.


        % \item Multimodal data from participants acting as a solid foundation for behavioral multimodal integration that supports a comprehensive analysis of sensorimotor and executive dysfunctions.
        % \item Offers a foundation for

        % using the VideoMAE-v2 model supports the computation of accurate and interpretable symbolic representations of the participants' egocentric vision stream, as described in \Cref{sec:chapter_5}. 
        
        
        % hand and skeleton pose landmarks using the MediaPipe library \cite{lugaresi2019mediapipeframeworkbuildingperception}, and speech and speaker diarization using the WhisperX model \cite{bain2022whisperx}. These features support a multimodal integration framework based on symbolic representations computation. 

        %  \item Features extraction include Extraction of speech and speaker diarization using WhisperX, hand and skeleton< pose landmarks using the MediaPipe software, and latent representations of video frames using VideoMAE-v2. These features support a multimodal integration framework based on symbolic representations computation.
        % \item Preparation (pre-processing, cleaning, registration, multimodal features extraction) of the Cooking Task (CT) and Arc de Triomphe Construction Task (ATCT) datasets comprised of 162 participants, 373 total administrations, with 1230 hours of multi-view videos (including egocentric video, gaze and oculometric data).

        % \item We extract speech and speaker diarization (WhisperX), hand and skeleton pose landmarks (Mediapipe software), and latent representations (VideoMAE-v2) of blocks of video frames to support a multimodal integration framework based on symbolic representations computation. 

    \end{enumerate}
\end{chapterabstract}


\section{Introduction}
\label{sec:chapter_4_introduction}

% 1) Presentation of neurological disorders in general, then executive function, and specifically TBI and LPR.
Executive Functions (EF) are high-level cognitive processes essential for planning, reasoning, decision making, inhibitory control, and goal-directed behavior \cite{Cristofori2019, Diamond2013}. These functions underpin problem-solving in complex, real-world situations and are fundamental to various aspects of daily life, including academic achievement, career success, social relationships, and public safety \cite{Robertson1997}, and were found in early childhood to be concurrently associated with emotion understanding and regulation, prosocial skills, peer acceptance \cite{Stucke2024}. Impairments in these functions, whether by injury, disease, or neurodegeneration, known as executive dysfunction or Dysexecutive Syndrome (DS), can severely affect an individual’s quality of life \cite{PabloRos2024, Rabinovici2015}. 
They are often linked to changes in the prefrontal cortex that come with normal aging \cite{Raz1999, Burke2006}, and are also frequently observed in people with neurological or psychiatric conditions involving damage to the frontal lobes \cite{Godefroy2003}. This includes individuals recovering from traumatic brain injury (TBI) or stroke \cite{Baddeley1988, Tsai2021}, as well as those living with chronic conditions such as Parkinson’s disease \cite{Aarsland2005} and multiple sclerosis \cite{Hajric2023}. For instance, it affects around half of individuals in the subacute phase following moderate-to-severe TBI \cite{Ribeiro2023}, underscoring its impact as a widespread and serious public health issue\footnote{See more information in  \cite{Borgnis2022} for a comprehensive systematic review of executive functions.}.




%Executive functions rely on a network of brain regions working together across the central nervous system. When these networks are disrupted - whether by injury, disease, or neuro-degeneration - DS can emerge \cite{Rabinovici2015}. 
% Specifically, these higher-order cognitive abilities and behavioral skills are responsible for controlling and regulating actions (e.g., starting and stopping activities or monitoring) (Burgess and Simons, 2005; Chan et al., 2008) and performing complex or non-routine tasks (e.g., ability to perform two tasks simultaneously) (Godefroy, 2003; Alvarez and Emory, 2006; Alderman, 2013).


%  These functions are important for many aspects of daily life, such as problem solving, decision making, and goal directed behavior. Dysexecutive syndrome affects the functioning of executive functions, and are closely related to difficulties in performing activities of daily living ~\cite{Diamond2013, Godefroy2003}, influences mental health, physical health, quality of life, school readiness, academic success, career success, marital life and public safety \cite{Diamond2013}.


% Impairments in EF—collectively referred to as Dysexecutive Syndrome (DS) or executive dysfunction—can severely compromise an individual’s autonomy and quality of life \cite{Robertson1997, Diamond2013, Cristofori2019, Godefroy2003}. DS is associated with difficulties in performing activities of daily living and has broad implications for both mental and physical health.


Diagnosing and quantifying DS is challenging due to the intricate nature of the high-order cognitive processes involved \cite{Tsai2021}. Dysexecutive syndromes are commonly assessed via standardized paper-and-pencil neuropsychological tests, typically administered in controlled clinical or laboratory settings \cite{Godefroy2010}. Current gold standard assessments include the Modified Wisconsin Card Sorting Test (WCST) \cite{Nelson1976}, the Trail Making Test~\cite{reitan1986trail}, and the Montreal Cognitive Assessment (MoCA) \cite{Freitas2013}, the Behavioral Assessment of Dysexecutive Syndrome \cite{Diamond2013}, and the Inventory of Dysexecutive syndromes (ISDC) \cite{Godefroy2010, Godefroy2008-sa}\footnote{For example, the Trail Making Test asks participants to alternate between connecting numbers and letters in sequence, assessing mental flexibility and divided attention. In the Wisconsin Card Sorting Test, participants must infer and adapt to changing sorting rules, thereby revealing their cognitive flexibility and ability to shift strategies.}.
While these tests demonstrate strong internal validity, they are limited in their ability to gauge the real-world impact of DS, known as their ”ecological validity”. In addition, these are often time-consuming, complex to administer, require trained clinicians or psychologists, and are subject to variability between examiners \cite{CHAN2008, Kessels2018, Borgnis2022}.
% Secondly, several studies revealed that many patients with Dysexecutive Syndrome achieve normal scores on traditional neuropsychological tests and, at the same time, complain of substantial difficulties in daily life activities (Shallice and Burgess, 1991).

% To address this, novel tasks and batteries have been developed to better reflect and disentangle everyday executive functioning, such as, the Behavioral Assessment of Dysexecutive Syndrome \cite{Diamond2013}, and the Inventory of Dysexecutive syndromes (ISDC) \cite{Godefroy2010, Godefroy2008-sa}. However, these are often time-consuming, complex to administer, require trained clinicians or psychologists, and are subject to variability between examiners \cite{CHAN2008, Kessels2018}.

% the Multiple Errands Test (MET) \cite{SHALLICE1991, Rand2009} 


%  However, these tests may not fully reflect an individual's real-life functioning, and may not capture the complexity and variability of dysexecutive symptoms in daily life.

To address this, researchers and clinicians have pursued the direct assessment of executive functions through tasks rooted in daily life. These include ecological tests, such as cooking or shopping activities \cite{Howieson2019, SHALLICE1991, Rand2009}, and semi-ecological tests, which mimic real-world cognitive demands without replicating daily routines \cite{Wilson1998, Cardoso2015,  Chevignard2008, Campos2023}. A test is considered ecological if it satisfies two criteria: (1) representativeness, meaning the task closely resembles a real-world situation in both form and content; and (2) generalizability, meaning that poor performance on the test reliably predicts functional difficulties in everyday life \cite{afdd1ed6d28448bb8c236ad8c5995a56}.

Building on these foundations, technological maturation have opened new avenues for assessing executive dysfunction with greater objectivity, granularity, and ecological relevance \cite{Bavaresco2022}. First, computational adaptations of traditional neuropsychological tests offer automatic stimulus delivery and precise measurement of response accuracy and timing, offering improved standardization and scalability \cite{JOHNSON2008, Camilleri2021}. Second, immersive technologies such as virtual reality (VR) and augmented reality, which aim to reproduce cognitively demanding situations in interactive realistic environments for evaluating complex cognitive functions, including planning, attention shifting, adaptation, multitasking, and problem-solving in scenarios that resemble everyday life \cite{Castelnuovo2003, Nolin2016, Parsons2018, LANIER201970, Realdon2019, Kourtesis2020, Borgnis2021, Borgnis2023}. Finally, the instrumentation of ecological and semi-ecological tasks using video recordings, eye-tracking and wearable sensors enable the collection of rich, multimodal behavioral data \cite{Campos2023}. 

While such technologies offer great potential for more accurate and ecologically valid assessments, their validation remains a critical step \cite{Renison2012}. In many cases, these approaches are still confined to research settings, and often rely on time-consuming procedures or subjective manual annotation of behavior.

The Smartflat for Dysexecutive Syndrome Stratification (SDS2) study aims to address this gap by identifying and evaluating behavioral and digital biomarkers of DS through a multidisciplinary approach combining traditional neuropsychological, ecological and semi-ecological tests, with multimodal data collection and specifically designed computational methods. The aim of the study is to develop a digital phenotyping approach to assess sensorimotor and executive dysfunctions across the following five dimensions: static balance, dynamic balance, and locomotion; neurovegetative and emotional state; fine motor skills; visual information processing; and planning, initiation, mental flexibility, and memory.

To this end, this part of the thesis focuses on the development of mathematical tools to quantify participant's behavior accurately and objectively. Our goal is to facilitate the study of common and specific aspects of behavior between healthy participants and those with acquired brain lesions. In this rest of the chapter, \Cref{sec:chapter_4_methods} provide a description of the study and data collection, and \Cref{sec:chapter_4_features_extraction} describes the features extraction methodology. Next,  \Cref{sec:chapter_5} describes our approach to estimate an interpretable vocabulary of egocentric vision prototypes, which is used to compute symbolic representations of the participants' vision stream. In \Cref{sec:chapter_6}, we develop novel adaptations to both the Time Warp Edit (TWE) distance and DTW Barycenter Averaging (DBA) methods for our symbolic sequences. Finally, we explore the ability of our computational framework to compare behavioral trajectories across stratified levels of dysexecutive syndromes severity. An overview of our agnostic approach to computational behavioral phenotyping approach is presented in \Cref{fig:chapter_4_overview_multimodal}.

%CODE: description of what we achieved
%A computational model then takes advantage of the observed fast temporal structure to describe mouse behavior as a series of reused modules, each a brief and stereotyped 3D trajectory through pose space that is connected in time to other modules through predictable transitions. 
 %Team members contributions are included in ~\Cref{app:contributions}.

%CODE: TODO IF TIME, the same miultimodal one but with the notion of sumbolic representation 

\newpage
\begin{figure}[!htbp]
    \centering
    \includegraphics[width = 0.9\linewidth]{chapters/chapter 4/figures/chapter_4_overview_multimodal.png}
    \caption{\textbf{Overview of our egocentric symbolic representation analysis framework.} (a) Multiple sensors record participants’ behavior during the cooking task, including three fixed cameras, an eye tracker capturing egocentric video, gaze, and oculometric data, as well as head movements and audio recordings. (b) Hand and body pose landmarks are estimated using the MediaPipe library \cite{lugaresi2019mediapipeframeworkbuildingperception}; speech transcription and speaker diarization are extracted with the WhisperX model \cite{bain2022whisperx}; and video segments of $\approx 640\,\text{ms}$ are encoded as sequences of latent representations using VideoMAE-v2 \cite{wang2023videomaev2}. (c) Egocentric vision prototyping and (d) unsupervised action boundary detection are used to obtain (e) fine-grained, interpretable symbolic representations of the participants’ egocentric video recordings. (f) We developed a novel elastic distance based on the Time Warp Edit (TWE) distance to compare symbolic sequences, used to(f) compute average typical representations of the cooking task by adapting the DTW Barycenter Averaging (DBA). (g) These are used to  investigate behavioral differences across diagnosis groups during the execution of the cooking task.}
    \label{fig:chapter_4_overview_multimodal}
\end{figure}


% BONUS/
% Therefore, it is recommended that the neuropsychological assessment include ecological tasks that provide similar scenarios to the patient`s daily life, being able to predict their performance in the real-world~\cite{Renison2012}. Studies after TBI highlight the importance of using tasks whose performance is evaluated through real situations that allow the professionals to observe the underlying difficulties~\cite{rand2018development}.


% The instrumentation of these tasks are often captured using wearable sensors, video recordings, or other digital traces, enabling rich, objective behavioral data to be collected \cite{Campos2023}. Despite their promise, these approaches remain relatively rare and are not yet widely adopted in clinical settings.


% Ecological and Computer Vision-Based Assessments
% Realdon2019 and Borgnis2022: Virtual Reality (VR) and other digital tools are increasingly used to provide ecological assessments of executive functions, simulating everyday activities in controlled yet realistic environments. VR-based tools enable dynamic interaction with 3D settings that replicate real-life scenarios, helping to capture executive dysfunctions as they manifest in daily life.
% Realdon2019: The PIT 360° tool is an example of an ecological measure using immersive technology to screen executive functioning in multiple sclerosis patients, capturing real-world executive dysfunction not always detected by conventional neuropsychological tests.
% Borgnis2022: Studies emphasize that traditional executive function tests often lack ecological validity and may not reliably predict functional impairments in daily life. Ecological assessments, including VR-based tasks and real-world simulations like the Multiple Errands Test (MET) or the Behavioral Assessment of the Dysexecutive Syndrome (BADS), provide more representative and generalizable insights into executive dysfunction.
% PabloRos2024, Borgnis2022 Research also reports that ecological assessments can detect executive impairments that traditional tests miss, and they are crucial for understanding how executive dysfunction affects everyday functioning, academic achievement, and behavior.
% Summary
% Camilleri2021 Recent research integrates machine learning mainly for analyzing psychometric and clinical report data related to executive dysfunction.
% Borgnis2022, Realdon2019 Ecological assessments using VR and immersive tools are recognized as essential for capturing executive dysfunction in real-life contexts, with some tools explicitly designed to simulate everyday tasks and environments.
% These ecological approaches are increasingly combined with advanced computational methods to improve diagnosis, monitoring, and rehabilitation of dysexecutive disorders.
% Thus, current trends strongly support the use of ecological assessment methods, often leveraging computer vision and VR technologies, alongside machine learning models to better understand and evaluate executive function impairments in real-world settings.

% or impairments 

\section{Study Description} 
\label{sec:chapter_4_methods}

The SDS2 study is a monocentric, longitudinal, non-randomized, COSMIN-compatible \cite{Terwee2018} diagnostic study aimed at comparing executive functions of brain-damaged patients with those of healthy individuals by assessing them with a multi-instrumented Cooking Task (CT) \cite{Chevignard2008}, and the Arc de Triomphe Construction Test (ATCT) \cite{Campos2023}. The study is conducted at the Hôpital National d’Instruction des Armées (HNIA) Percy, France, by the neurology team of the hospital. The recruitment of participants is described in \Cref{apd:chapter_4_methods}.

The Cooking Task used in this study is an ecological test based on Chevignard’s work\cite{Chevignard2008}, and consists in cooking a chocolate cake according to mandated instructions from a manual. Previous findings indicate that cooking can provide a sensitive and reliable measure of executive function ability in a “real-world” context \cite{Chevignard2008, Doherty2015}. The act of cooking a meal requires several executive functions, including capacity to multitask, plan, use prospective memory, and maintain and complete both sub and overall goals within a strict timeframe \cite{craik2006cognition}. In practice, patients with DS tend to make more errors, such as context neglect, purposeless action, or poor environmental adherence. We report in \Cref{apd:chapter_4_error_type_description} a description of the different types of errors that can be observed during the cooking task.

The ATCT task was designed at the Borelli Center to be highly specific to reveal DS, and consists in a one-hour test that consists in building a model of the Arc de Triomphe with 541 Lego$¨\circledR$ bricks in 131 separate steps. It was proposed as a highly discriminant task to evaluate DS due to its length and its design. 

%It extends previous works led by the same team, where these tests had been  administered to participants meeting similar inclusion criteria. Data from these par-  ticipants do not belong to the SDS2 study and will therefore not be included in the  corresponding statistical analysis. However, diagnostic thresholds determined in this  population will be used when assessing the performance of any test under investigation  to mitigate the risk of overestimation bias [10].  



While the methodology introduced in this part of the thesis is applicable to data collected from this task, and is referenced throughout the manuscript, \textbf{the present analysis focuses exclusively on the cooking task.} The application of our approach to the ATCT task is left for future work.





% We develop a multi-instrumented cooking test, the Smartflat study, which is administered in clinics at the neurology department of the Military Hospital of Percy (France), and presents a cooking task, a Lego task while the patient's behavioral and oculometric responses are recorded via three fixed cameras and wearable eye-tracking glasses recording egocentric vision. 



% The data collected consisted of video recorded by three cameras placed at different viewpoints of the room, as well as an oculometer (or eye-tracker) worn by the participant, which was comprised of embedded accelerometers and gyroscopes.


\subsection{Cooking task administration and data collection}
\label{sec:chapter_4_app_preesentation}

To perform this test, each participant is asked to stand at a fully-equipped kitchen where the cookbook presents the recipe and necessary ingredients are disposed at specific fixed positions (see \Cref{fig:chapter_4_material}). The test ends when the participant considers the task complete or is unable or unwilling to continue.

\begin{figure}[!htbp]
    \centering
    \includegraphics[width = 0.6\linewidth]{chapters/chapter 4/figures/chapter_4_material.png}
    \caption{\textbf{Initial setting of the Smartflat kitchen before performing the Cooking Task.} While reproducibility constraints are crucial for the development of standardized assessment tests, the results presented later in \Cref{sec:chapter_5} reinforce the importance of making the test environment visually reproducible (e.g., same ingredients and ustensils).}
    \label{fig:chapter_4_material}
\end{figure}

The task is instrumented using three GoPro Hero 5 cameras, strategically located to record at a sampling frequency of 60 frames per second the participant and their environment, as illustrated in \Cref{fig:chapter_4_CT_thumbnails}.

\begin{figure}[!htbp]
    \centering
    \includegraphics[width = 1\linewidth]{chapters/chapter 4/figures/chapter_4_CT_thumbnails.png}
    \caption{\textbf{Representative samples showing the four cameras angles.}}
    \label{fig:chapter_4_CT_thumbnails}
\end{figure}

Before the test starts, the participant is equipped with a Tobii Pro Glasses 2 portable eye-tracker glasses, equipped with an embedded microphone, an ego-centric camera and an inertial measurement unit composed of three gyroscopes and three accelerometer recording head movements in each direction. These glasses record the participant's eye movements, gaze patterns, and head movements at a sampling frequency of 100Hz, while the embedded camera captures egocentric video at a resolution of 1920x1080 pixels and a frame rate of 25 frames per second. These glasses provide a first-person perspective of the participant's actions and interactions with the environment, capturing their visual attention and spatial orientation during the task. The combination of the cameras and oculometric glasses allows for a detailed analysis of the participant's behavior, providing insights into their cognitive and motor functions.

Trained examiner first explains the task to the participant, ensuring they understand the instructions and what is expected of them. The examiner emphasizes the importance of following the recipe closely and completing the task to the best of their ability. The participant is then asked to prepare a chocolate cake by following the provided recipe, which includes a list of ingredients and step-by-step instructions. An examiner remains present in the room to ensure the participant’s safety and to provide timely assistance if the participant is clearly unable to proceed to the next step.

% Three Go Pro Hero H cameras record the participant from different angles, providing a comprehensive view of their behavior throughout the task. A
Before and after the assessment, participants complete questionnaires assessing their emotional state and their level of proficiency with cooking and Lego construction tasks.
Lastly, a trained neuropsychologist reviews the recorded videos to manually annotate the starting and ending times of the test, ensuring the analyzed data corresponds accurately to the task duration only.


\subsection{Analyzed cohort}
\label{sec:chapter_4_analyzed_cohort}  
150 participants between 19 and 85 years old were enrolled in the study starting in February 2017, with a total of 175 administrations. At the end, 143 participants and 162 administrations were included in the analysis (15 participants performed the task twice, and 2 three times) if data collection met minimal quality criteria (complete and manually validated videos and eye-tracking data). Administrations were excluded if the egocentric video were missing. We report available demographics, clinical variables, and total test duration in \Cref{table:chapter_4_demographics_clinical}. % CODE: The average task duration

%The average number of administrations per participants for both the lego and cooking tasks are 1.5 (std=X) and 2.1 (std=Y) respectively. Modalities were removed if under the minimal quality requirements (see section X). For the lego task (N=132), this resulted in X participants having their 4 modalities passing  quality checks, X with 3 modalitities, X with 2 modalities, and X administrations with only one modalitiies. See the empirical histograms of missing data per task in Figure X.

% Quality criteria:  Not corrupted (able to collate, meeting temporal consistency criteria (i.e. all non-corrupted modalities must have a maximum total duration under e.g 10 minutes. (Show histogram ?)



% \begin{table}[!htbp]
% \centering
% \begin{tabular}{p{4.0cm}cccc}
% \hline
% \textbf{Sample Desc.} & \textbf{Control (N=19)} & \textbf{TBI (N=51)} & \textbf{RIL (N=35)} & \textbf{OTH (N=9)} \\
% \hline
% \textbf{Age Mean (SD)} & 33.0 (7.8) & 36.8 (11.5) & 53.6 (16.5) & 50.7 (15.9) \\
% Unknown Total (\%) & 16 (84.2) & 2 (3.9) & 3 (8.6) & 0 (0.0) \\
% \textbf{Male Total (\%)} & 0 (0.0) & 41 (80.4) & 17 (48.6) & 4 (44.4) \\
% \textbf{Female Total (\%)} & 19 (100.0) & 10 (19.6) & 18 (51.4) & 5 (55.6) \\
% Unknown Total (\%) & 0 (0.0) & 0 (0.0) & 0 (0.0) & 0 (0.0) \\
% \textbf{ISDC Mean (SD)} & -- & 4.7 (14.4) & 4.2 (8.2) & 2.0 (3.5) \\
% Unknown Total (\%) & 19 (100.0) & 16 (31.4) & 13 (37.1) & 4 (44.4) \\
% \textbf{MoCA Mean (SD)} & -- & 24.7 (4.7) & 21.2 (4.2) & 22.7 (2.5) \\
% Unknown Total (\%) & 19 (100.0) & 18 (35.3) & 14 (40.0) & 2 (22.2) \\
% \textbf{Duration Mean (SD)} & -- & 24.7 (4.7) & 21.2 (4.2) & 22.7 (2.5) \\
% Unknown Total (\%) & 19 (100.0) & 18 (35.3) & 14 (40.0) & 2 (22.2) \\
% \hline
% \end{tabular}
% \caption{\textbf{Demographic and Clinical Characteristics of Study Sample.} TBI: traumatic Brain Injury; RIL: Radiation-Induced Leukoencephalopathy. OTH: Other neurological pathology. ISDC: Inventaire du Syndrome Dysexécutif Comportemental \cite{ISDC, Godefroy2008-sa}. MoCA: Montreal Cognitive Assessment \cite{MoCA}. Unfortunately, 29 participants have missing demographics and clinical data and are thus exluded from the table. These statistics aim to be updated as to include all 143 participants as soon as available.}
% \label{table:chapter_4_demographics_clinical}
% \end{table}


\begin{table}[H]
    \centering
    \begin{tabular}{p{4cm}cccc}
    \hline
    \textbf{Groups} & \textbf{Control (N=26)} & \textbf{TBI (N=59)} & \textbf{RIL (N=37)} & \textbf{OTH (N=5)} \\
    \hline
    \textbf{Age Mean (SD)} & 40.31 (10.06) & 38.66 (14.19) & 53.69 (15.47) & 52.60 (16.92) \\
    \quad Unknown Total (\%) & 10 (38.5\%) & 12 (20.3\%) & 8 (21.6\%) & 0 (0.0\%) \\
    \textbf{Male Total (\%)} & 3 (12.0\%) & 43 (82.7\%) & 20 (58.8\%) & 1 (20.0\%) \\
    \textbf{Female Total (\%)} & 22 (88.0\%) & 9 (17.3\%) & 14 (41.2\%) & 4 (80.0\%) \\
    \quad Unknown Total (\%)  & 1 (3.8\%) & 7 (11.9\%) & 3 (8.1\%) & 0 (0.0\%) \\
    \textbf{ISDC Mean (SD)} & -- & 5.00 (15.01) & 3.85 (8.10) & 2.67 (4.62) \\
    \quad Unknown Total (\%) & 26 (100.0\%) & 27 (45.8\%) & 17 (45.9\%) & 2 (40.0\%) \\
    \textbf{MoCA Mean (SD)} & -- & 24.74 (4.73) & 21.55 (4.05) & 22.50 (3.42) \\
    \quad Unknown Total (\%) & 26 (100.0\%) & 28 (47.5\%) & 17 (45.9\%) & 1 (20.0\%) \\
    \textbf{Duration Mean (SD)} & 28.21 (6.11) & 35.39 (10.72) & 42.07 (9.54) & 42.45 (8.57) \\
    \hline
    \end{tabular}
    \caption{\textbf{Demographic and Clinical Characteristics of Study Sample.} TBI: traumatic Brain Injury; RIL: Radiation-Induced Leukoencephalopathy. OTH: Other neurological pathology. ISDC: Inventaire du Syndrome Dysexécutif Comportemental \cite{ISDC, Godefroy2008-sa}. MoCA: Montreal Cognitive Assessment \cite{MoCA}. Unfortunately, 29 participants have missing demographics and clinical data and are thus excluded from the table. These statistics aim to be updated as to include all 143 participants as soon as available.}
    \label{table:chapter_4_demographics_clinical}
\end{table}





% \subsection{Inclusive and ethical conduct of research}

% Consent was obtained from all individuals whose faces are shown in the figures or videos for publication of these images. Participants provided written informed consent, and the study protocols were approved by the Comité de Protection des Personnes (CPP) Sud Mediteranee IV (CPP: 210504). Trial is registratered at \url{ClinicalTrials.gov} with identifier: NCT05017051, IDRCB: 2021-A00087-34. The algorithms presented in this work are described therein to enhance reproducibility, or available in their official repositiories. For additional access or request, please contact the corresponding author at \url{sam.perochon@ens-paris-saclay.fr}.

% \paragraph{Data availability.\\}

% In line with the Open Science principles [21], the reproducibility policy of the Centre Borelli, regulatory requirements, and technical feasibility, anonymized and full or partial data will be made available upon request, which include the raw eye-tracking data, anonymized versions of the recorded videos, and the intermediary algorithms outputs. 

% To ensure the full transparency of our work and stimulate collaborations and secondary analyses of the collected data, and regulatory
%  we aim to publicly release as much collected data as possible. Depending on regulatory issues at the time of publication, this may include releasing raw eye-tracking data as well as anonymized versions of the recorded videos, and intermediary relevant outputs produced by the algorithms developed. If possible, algorithms developed as part of this study will be deposited in open source repositories. To ensure the full transparency of our work and stimulate collaborations and secondary analyses of the collected data, we also plan to release a fully synthetic version of our dataset generated using a non-parametric method recently developed by Chapelle et al [22].  


\section{Features extraction}
\label{sec:chapter_4_features_extraction}

The Smartflat dataset comprises multi-view videos, gaze and oculometric data, head kinematics, and audio recordings, capturing diverse aspects of participant behavior. \Cref{sec:chapter_5} presents a method to estimate fine-grained and interpretable symbolic representations of participants’ egocentric vision recordings during the task. Our aim is to study how lower-dimensional features extracted from these additional modalities can be integrated into the egocentric vision–based symbolic representation. This section focuses on the extraction of primary features of interest derived from these modalities. The dataset preparation phase of the thesis involved encoding blocks of video frames into latent representations using a video foundation model, estimating hand and skeleton landmarks from two camera views, transcribing speech and identifying speakers, analyzing gaze patterns. An overview of our multimodal output extraction pipeline is shown in \Cref{fig:chapter_4_overview_multimodal}-a.\footnote{This work was performed using computational resources from the “Mésocentre” computing center of Université Paris-Saclay, CentraleSupélec and École Normale Supérieure Paris-Saclay supported by CNRS and Région Île-de-France ([https://mesocentre.universite-paris-saclay.fr/](https://mesocentre.universite-paris-saclay.fr/)). 
}


% \begin{figure}[!htbp]
%     \centering
%     \includegraphics[width = 0.9\linewidth]{chapters/chapter 4/figures/chapter_4_overview_multimodal.png}
%     \caption{\textbf{Multimodal outputs extraction from the Smartflat dataset.} The figure illustrates the different modalities and their respective outputs extracted from the Smartflat dataset. The modalities include multi-view video frames, gaze and oculometric data, head movements, and audio recordings. These unimodal streams are processed to obtain hand and pose landmarks, gaze events estimation, video blocks latent embedding sequences, and transcripted speech and speaker diarization.}
%     \label{fig:chapter_4_multimodal_outputs}
% \end{figure}

While a detailed analysis of landmark estimates, gaze, oculometric signals, and speech patterns is left as future work, our objective is to present our feature extraction methodology, provide illustrative examples, and explore how these modalities can support the assessment of dysexecutive syndromes in the context of the Smartflat study. In particular, in \Cref{sec:chapter_6_multimodal_integration}, we discuss how the egocentric visual context encoded in interpretable symbolic representations can facilitate multimodal analysis, and support a more systematic and effective integration of diverse modalities.


\subsection{Encoding block-of-frames into compact latent representations using video foundation models}


% Shall we integrate in filigram here that the search for local descriptors 

% Here I want to mention in a subtle way the challenge of image and video analaysis by paying a tribute to the old fashion algorithms. 
% Before the emergence of large-scale datasets and successful deep neural networks architecture and learning objectives, image and video processing algorithms focused in keypoints extractions, the design of local descriptors (e.g.  Fisher vectors, SIFT, ORB), and used optical flow computation to encode movements. (Here I want to achnoledge more of the variety and creativity researchers demonstrated in the field, but remain really really "down-to-earth' as it is the case here).

% Here I want to introduce foundation models, and then decline for video data. 
Inspired by their success in natural language processing \cite{Bert, openai2024gpt4technicalreport}, large-scale training of foundation models via self-supervised learning (SSL) have demonstrated the ability to learn rich, compact, and transferable representations across modalities — including images \cite{Radford2021LearningTV}, audio \cite{pmlr-v162-baevski22a}, and physiological time series \cite{Banville2021-qz, abbaspourazad2024largescale}.
% Start with Our study suggests that the general framework of masked autoencoding (BERT [15], MAE [31], etc.) can be a unified methodology for representation learning with minimal domain knowledge.


%These representations are learned without explicit labels, by solving pretext tasks designed to capture structure inherent in the data.

% Following their success in natural language processing \cite{Bert, openai2024gpt4technicalreport}, large-scale training of "foundation" models via self-supervised learning (SSL) demonstrated remarkable ability to learn rich and compact representations of image \cite{Radford2021LearningTV}, sound \cite{pmlr-v162-baevski22a}, and time series \cite{Banville2021, abbaspourazad2024largescale}.

Compared to images or text, videos typically exhibit a lower information rate, due to both high spatial redundancy across pixels and strong temporal correlation between adjacent frames \cite{videomae}. This encourages the computation of lower-dimensional general-purpose representations that can address challenging downstream tasks related to video understanding, such as action localization, recognition, anticipation,  question answering, captioning, or summarization \cite{ego4D}, requiring different degrees of abstract, temporal, and spatio-temporal understanding of video content \cite{madan2024foundationmodelsvideounderstanding}. 
%require motion understanding or long-term temporal reasoning

% Citing here the review for the first time before introducing it explicitly later in the text is a bit annoying; you could potentially change that  
A recent survey of video foundation models (VFM) identifies three categories of models: (i) image-based, which adapt frame-level image encoders to videos \cite{resnet}; (ii) video-based, pre-trained on video slices directly \cite{videomae}, and (iii) multimodal, which integrate video with additional modalities such as audio, text, or static images during pretraining \cite{lin2022egocentricvideolanguagepretraining, Desai2023, li2024videovistaversatilebenchmarkvideo, madan2024foundationmodelsvideounderstanding}. Preliminary VFM benchmark studies tend to show that image-based FMs performs better on tasks requiring accurate visual encoding, while video-based FMs tend to achieve better on tasks requiring motion understanding or long-term temporal reasoning \cite{madan2024foundationmodelsvideounderstanding}

% Enorme question de travail de precision sur l'apprentissage par constrastrivité

The majority of SSL strategies fall into two paradigms, inducing different but present inductive biases in the learnt representation:

\begin{itemize}
    \item Contrastive learning, as exemplified by CPC \cite{Oord_Oriol_2018}, SimCLR \cite{pmlr-v119-chen20j}, DINOv2 \cite{oquab2023dinov2}, and others, where the goal is to discriminate between similar and dissimilar samples in a learned latent space; and
	\item Generative or reconstruction-based learning, such as masked autoencoding (MAE) \cite{MAE, videomae}, where the model is trained to reconstruct masked portions of the input from visible patches, promoting compactness and information compression.
\end{itemize}

Contrastive methods tend to produce representations that emphasize semantic separation and depend on the chosen definition of contrastivity, while masked reconstruction objectives are generally considered more agnostic, as they learn to capture the internal data structure without explicitly relying on instance discrimination.

% While contrastive pre-training (name a few model names) have been proposed as efficient VFM, they often rely on the chosen definition of contrastivity, compared to reconstruction-based strategies that are more agnostic. These two learning mechanisms induce different learning biases at the representation level, which may be desirable or not depending on the application. 
Since our goal is to encode videos as a pre-processing step rather than to address a predefined video-understanding task, we adopt a reconstruction-based approach, more specifically the VideoMAE-v2 architecture — a state-of-the-art masked autoencoder for video representation learning. VideoMAE-v2 builds on the success of masked autoencoding in the image domain by extending it to spatiotemporal data \cite{MAE, videomae-model}. Both encoder and decoder use transformer building blocks that jointly models spatial and temporal attention \cite{dosovitskiy2021imageworth16x16words}, and are trained with an optimal masking ratio of 90\% of the spatiotemporal "tubes" - which is related to the high level of information redundancy in video data \cite{videomae-model, wang2023videomaev2}.

\paragraph{Embedding sequences extraction procedure.\\}

Let $V = {v_1, v_2, …, v_N}$ denote a sequence of $N$ video frames sampled at \textsc{fps} frames per second. We partition $V$ into overlapping windows of $w=16$ frames, with a stride of $w=8$ frames — resulting in 50\% overlap between spatiotemporal patches. Each blocks of frames $w_i=(v_i, ..., v_{i+w})$ is passed through the \texttt{giant} version of the pre-trained VideoMAE-v2 model (1.012 billion 16-bit parameters), which outputs a $D$-dimensional latent embedding $x_i \in \mathbb{R}^D$, with $D=1408$. The full video sequence is thus encoded as:

$$X = \{x_1, x_2, …, x_T\}, \quad x_i \in \mathbb{R}^D, \quad T = \left\lfloor \frac{N - 16}{8} + 1 \right\rfloor.$$

At 25 fps for the egocentric video, a 16-frame window covers 640ms, and a stride of 8 frames produces a new embedding every 320ms. This yields around 1k embedding every 5 minutes of videos, and an average total sequence length of 5592 (STD=1766) embeddings per participant. A total of $N=\sum_{i=1}^{N_s}N_i = 906905$ embeddings encode the egocentric video of $N_s=162$ subjects having each $N_i$ embeddings representing their visual experience during the cooking task. These latent sequences constitutes the primary material for downstream behavioral analysis methods, such as the egocentric vision prototyping presented in \Cref{sec:chapter_5}. 

\subsubsection{Preliminary analysis of egocentric vision latent sequences}
\label{sec:chapter_4_preliminary_analysis_egocentric_latent_sequences}
% This hierarchical structure makes the definition of an exhaustive vocabulary of high-level action labels impractical.  
Typically, the comparison of encoded samples is quantified using the cosine similarity between their representations. Given two vectors $u, v \in \mathbb{R}^D$, their cosine similarity is defined as $\text{sim}(u,v) \;=\; \frac{u \cdot v}{\lVert u \rVert \,\lVert v \rVert},$  and measures the angle between $u$ and $v$ in the representation space. We note $d$ their cosine distance, defined as $d(u,v) \;=\; 1 - \text{sim}(u,v)$. 

In theory, the hierarchical structure of images — from low-level cues (e.g., edges, corners, orientations), mid-level features (e.g., shapes, textures), to high-level representations (e.g., objects, scenes, semantics, etc.) — implies that contiguous blocks of frames of the same action are expected to present similar appearance, and thus a larger pairwise similarity. In contrast, different actions are expected to exhibit distinct visual signatures, whose degrees of visual differences potentiate the ability for similarity-based approach to classify them as different actions. To examine this, we compute the pairwise cosine similarity between all time-aligned embeddings of a participant’s egocentric video, called Gram matrix, defined as $G_{i,j} = \tfrac{x_i^Tx_j}{|x_i||x_j|}, \quad (i, j) \in [N]^2 $, where $x_i$ denotes the latent representation at time step $i$. We visualize in \Cref{fig:gram} the Gram matrix of a short sequence of latent representations, corresponding illustrative frames, and the ground-truth actions annotations.  % We report in the following a qualitative analysis of 5 minutes of a sequence egovision embeddings to support our methodology in \Cref{sec:chapter_5} and \Cref{sec:chapter_6_barycenter_estimation}. % To support the development of our methodology in \Cref{sec:chapter_5} and \Cref{sec:chapter_6_barycenter_estimation}, we present here some key observations drawn from a preliminary analysis on the latent embedding sequences of the egocentric video.  
 % and draw the following key observations supporting the development of our methodology in \Cref{sec:chapter_5} and \Cref{sec:chapter_6_barycenter_estimation}. 

\begin{figure}[!htbp]
\centering
\includegraphics[width = 0.9\linewidth]{chapters/chapter 4/figures/gram_photo.png}
\caption{\textbf{Pairwise cosine similarity matrix between time-aligned embeddings, representing 3 minutes of egocentric vision content for a randomly sampled participant.} The ground truth (GT) depicts manual action segmentation collected by neurologists.}
\label{fig:gram}
\end{figure}

From this illustrative example, we define three structural patterns in the Gram matrix:
\begin{enumerate}
    \item \textbf{Diagonal blocks} — squared regions of high self-similarity along the diagonal reflecting stable egocentric vision episodes. At the lowest scale, they typically correspond to our atomic actions, such as reading, beating eggs or pouring ingredients, where short visual inputs remain consistent over several seconds. At larger scale, they can be mapped to actions with higher-level semantics, within which visual similarity is more similar within-group than across neighboring regions. 

    \item \textbf{Off-diagonal blocks} — regions exhibiting high pairwise similarity with their concurrent diagonal blocks, indicating repeated actions (e.g., cracking multiple eggs or reading the recipe). We propose in the next chapter to cluster these re-occurring patterns together to promote low-dimensional representation of behavior, considering the analysis of repeated actions and participant's progression through the task as post-hoc symbolic sequences analysis. 

    \item \textbf{Cross-shaped artifacts} — thin vertical and horizontal low-similarity bands often emerge from brief intervals of egocentric motion blur or noisy transitions (e.g., rapid head movements, occlusions, or abrupt viewpoint changes). These episodes produce poorly encoded latent representations that are uniformly dissimilar to the rest of the embedding sequences. These are typically encoded with a unique "background" class to prevent their ambiguity from impacting symbolic representations of participant's behavior. % We examine their effect empirically in \Cref{sec:chapter_5_noise_encoding}, where we observe that except for some long-enough or frequent episodes of this type that may integrate the vocabulary, they tend to be encoded as noisy segments to improve the quality of the ultimate symbolic representations

\end{enumerate}

Two central questions arise for the analysis of embedding sequences: (i) how to detect transitions between action units, and (ii) how to organize these units into coherent categories. We aim to address these challenges through two complementary tasks:

\begin{enumerate}
    \item \textbf{Unsupervised temporal action segmentation (uTAS)} (\Cref{sec:chapter_5_unsupervised_temporal_segmentation}), which identifies meaningful transitions in the embedding dynamics.
    \item \textbf{Unsupervised action clustering (uAC)} (\Cref{sec:chapter_5_prototyping}), which groups egocentric vision atoms into semantically consistent categories.  
\end{enumerate}

Together, uTAS and uAC can be viewed as unsupervised analogues of the classical temporal action localization (TAL) and action recognition tasks. While recent progress in large-scale self-supervised video representation learning and vision-language foundation models has reshaped how action understanding is approached, these methods still rely heavily on indirect supervision that is scarce in clinical research (e.g., large-scale pretraining on massive curated datasets or alignment with textual descriptions). By contrast, our similarity-based approach is deliberately agnostic to predefined taxonomies: it operates directly on the dynamics of video latent sequences to extract fine-grained action units. This unsupervised formulation is particularly valuable in contexts where annotations are scarce, costly, or ill-defined, and where discovering the fine-grained action units is needed to recover subtle behavioral patterns. In this sense, our framework offers a principled step toward transforming egocentric video into symbolic representations tailored for downstream behavioral and clinical analysis.



%  A key contribution of our work is to propose a data-driven prototype-based method to infer such hierarchical behavioral inventory from egocentric visual data collected during a cooking task (presented in \Cref{sec:chapter_5}).


% resources: lter work on commputational ethogram computation: We adopt this bottom-up approach in the next chapter, in which we ai and investigate how fine-grained atomic units of action, which we call egocentric vision atoms, can be used to infer actions of higher-level semantics and larger temporal scales.




% Ressource: link between visual hierarchical similarity and hierarchical semantic structure of human activities 
% We hypothesize that if the video encoder effectively encode this hierarchical structure, if we assume that vision characterize action semantic, and if human activity present such a hierarchical structure \cite{Rosenbaum2007, lashley1951serial}, then the visualization of time-ordered cross-similarity map are expected to reveal the organization of behavior as an hybrid between (i) sequential atomic actions at low granularity, and (ii) a hierarchical structure emerging at larger temporal scale. 

% \item \textbf{Intermediate regions} — the residual similarity structure, excluding the above three types, encodes the varying degrees of dissimilarity between different action segments. The dispersion of similarity values offer insight into the video encoder’s ability to define a meaningful metric space over egocentric visual domain. A well-structured latent space should assign low similarity to semantically unrelated segments, while preserving higher-order relationships (for exampled grasping similar objects, washing ustensils, etc.). Therefore, it may be that "grasping" actions of different objects may appear altogether more similar than breaking the eggs and then throwing the shells. 


% While happening in sequential order, action structure were shown to be hierarchical in their semantic structure, as illustrated in 

% \section{Conclusion}

%     Dataset: Standardization des datasets Lego + Gateau + procédures + documentations 
%         Stade III:  "analyzable" ou "publiable", xxx

%     % Turn to algo ? 
%     \paragraph{Steps of the procedure:\\}

%     \begin{enumerate}
        %COODE incorporate ? 
%         \item $X\leftarrow f_\theta(V_{3D})$ Estimation of the egocentric spatiotemporal tokens (Representation Learning)
%         \item $T_e, \hat{\lambda_0} \leftarrow Pelt(X)$ or $T_e = TwFinch(X)$ Temporal segmentation of the spatiotemporal grid (remove temporal correlation and keep information content)
%         \item $\tau_T, \tau_E, \tau_N \leftarrow  KernelDensityEstimation(S_e, \mu_T, \mu_E, \mu_N)$ Empirical modeling of the temporal distributions, with modes-decomposition
%         \item $S_e, \mu_T, \mu_E, \mu_N \leftarrow  Symbolization(X, T_e)$ Symbolization of the task and key spatiotemporal tokens extraction
%         \item $\overline{S_{control}} = ShapeTimeWassersteinDBA(S_e, \tau_T, \tau_E, \tau_N)$ Barycenter averaging (potentially several)
%         \item $\pi_{P \rightarrow C}, D_{\pi_{P \rightarrow C}} = (D_{ij})_{i, j \in \pi_{P \rightarrow C}} \leftarrow  ShapeTimeWassersteinDTW(S_e^P, \overline{S_{control}}, \tau)$ Temporal localization of errors
%     \end{enumerate}

%     The step  $S_e, \mu_T, \mu_E, \mu_N \leftarrow  Symbolization(X, T_e)$, includes the refinement recursive procedure, in order first to explore the possible visual tokens, before merging them based on sillhouette-score optimization.  Co-clustering at this step using the empirical kernel density estimates (KDE) of the temporal distributions of the symbols should be beneficial.


% Otherwise they will be put next to each other in the final barycenter estimate (maybe), and so it just increases artificially the amount 

        
% Describe the links between components and tasks ? 
% In this work, we hypothesized that in egocentric videos, actions should appear as a succession of frames with similar visual appearance (e.g., containing the same objects). Therefore, we expect features extracted from the images to have high intra-segment similarities and low inter-segment similarities. Conversely, in egocentric videos, the background segments are expected to be composed of a succession of images with low inter-segment similarities (e.g., due to head or body movements resulting in motion blur in the frames).




% Importantly, extensive literature on hierarchical visual processing in neuroscience and perception studies \cite{Kruger2013}, and is associated with the hierarchical nature of visual content (with have been extensively studied for image processing and deep learning. We therefore expect the representations of temporally close blocks to exhibit similar hierarchical structures.

% Lastly, the hierarchical nature of visual content, from low-level features (e.g. edges, orientations) to mid-level features (e.g. shapes, textures) to high-level representations (e.g. objects, scenes, semantics, etc.), allows to hypothesize that representations of temporally close blocks to exhibit similar hierarchical structures quantified by the choice of a metric (e.g. cosine, euclidean).%, and acknowledging the role of movement in modulating 
% This is illustrated in \Cref{fig:gram}, in particular during the the red and green segments, which 

% In particular, it helps assess the encoder’s ability to map temporally and semantically related events into well-structured and expressive regions of the latent space.




%What may allow to decipher between the use of the transition matrices or the latent space /temporal space distance use for ethogram inference may lies where the notion of hierarchical nature of behaviors and hierarchical nature of visual content cannot readily be disentangled and are not correlated 100%.


% What is the correlation between our -1 encoding and the average rank of embedding noise . Here we depart from a full use of the participant's data for denoising puprose and to enhance the quality of our symbolic representations. 



% \paragraph{Multimodal foundation models.\\}

% Recently, multimodal foundation models have been proposed to jointly model different modalities, such as text, images, and videos. These models leverage large-scale pre-training on diverse datasets to learn rich representations that can be fine-tuned for specific tasks. 


% \cite{Ashutosh2023}  Video-language embeddings are a promising avenue for injecting semantics into visual representations, but existing methods capture only short-term associations between seconds-long video clips and their accompanying text. We propose HierVL, a novel hierarchical video-language embedding that simultaneously accounts for both long-term and short-term associations. As training data, we take videos accompanied by timesta∏mped text descriptions of human actions, together with a high-level text summary of the activity throughout the long video (as are available in Ego4D). We introduce a hierarchical contrastive training objective that encourages text-visual alignment at both the clip level and video level. While the clip-level constraints use the step-by-step descriptions to capture what is happening in that instant, the video-level constraints use the summary text to capture why it is happening, i.e., the broader context for the activity and the intent of the actor. Our hierarchical scheme yields a clip representation that outperforms its single-level counterpart as well as a long-term video representation that achieves SotA results on tasks requiring long-term video modeling.

% \cite{lin2022egocentricvideolanguagepretraining} Egocentric Video-Language Pretraining


% \cite{li2024videovistaversatilebenchmarkvideo} VideoVista: A Versatile Benchmark for Video Understanding and Reasoning


% \cite{Desai2023} Hyperbolic image-text representations. Visual and linguistic concepts naturally organize themselves in a hierarchy, where a textual concept "dog" entails all images that contain dogs. Despite being intuitive, current large-scale vision and language models such as CLIP (Radford et al., 2021) do not explicitly capture such hierarchy. We propose MERU, a contrastive model that yields hyperbolic representations of images and text. Hyperbolic spaces have suitable geometric properties to embed tree-like data, so MERU can better capture the underlying hierarchy in image-text datasets.


% In the context of the Smartflat study, we can explore the use of such models to integrate the multimodal data collected from the cooking task, including video frames, eye movements, and speech data. By leveraging these models, we can potentially improve the accuracy and robustness of our behavioral analysis and digital phenotyping approach.



% Each modality stream were processed during the dataset preparation stage of the thesis to obtain hand and pose landmarks, gaze events estimation,  transcripted speech and speaker diarization, and video blocks latent embedding sequences. 
 

% We describe in the rest of this chapter our features extraction methodology, investigating each modality independently, and then combining the outputs to obtain a multimodal representation of the participants' behavior. 
% While we believe such integration cannot happen before each modality  

% Here draw a schematic showing the different ways (in an holistic fashion) to see modality streams, features extraction, and multimodal integration.


% \subsection{Eye movements and gaze estimation}
% \cite{Kempa2021} (in-house) Apports de l’oculométrie dans l’analyse d’une t\^ache procédurale
% \cite{lai2024listen} Listen to Look into the Future: Audio-Visual Egocentric Gaze Anticipation


% Vision Toolkit features ? 


% \subsection{Speech recognition and speaker diarization}

% \cite{rvad_fast}


%  \begin{figure}[!htbp]
%     \centering
%     \includegraphics[width=.6\linewidth]{chapters/chapter 2/figures/bubble_game.png}
%     \caption{\textbf{Frame of the Bubble Popping Game.} This game is composed of 5 Tracks with bubbles appearing from the bottom Any time a bubble is popped by touching it with a finger (whether by caregiver or child), the bubble pops with a distinct popping sound and the animal is now fully seen and spins around while little fireworks explode around the animal. The animal floats off the screen while other bubbles with animals in them continue to float by.}
%     \label{fig:bubble_game}
%  \end{figure}


\subsection{Gaze, oculometric, and pupillometric features extraction}
\label{sec:chapter_4_eye_tracking}
     % paper: add inhouse conclusions 10.1145/3486812.3486841
    Eye-tracking and pupillometry have emerged as powerful tools for investigating cognitive load and arousal, attentional control deficits, sensorimotor anticipation, and visual-motor control deficiencies in both healthy individuals and clinical populations \cite{Wolf2023, Opwonya2024, Tao2020}. Eye movement metrics, such as saccade dynamics (e.g., latency, velocity, spread), fixation patterns (durations, frequency), and pupil responses (diameter dynamics, light reflex, task-evoked dilation \cite{Castellotti2025, Petrosino2025, Strauch2022}) were shown to provide sensitive biomarkers of mild cognitive impairments (MCI) and executive functioning \cite{Graham2024,Tao2020}. However, traditional paradigms often rely on constrained laboratory tasks (e.g., antisaccade, n-back) that may not fully capture the complexity of real-world executive functioning. Recent advances in mobile eye-tracking technology now enable the collection of gaze and pupillometry data in more naturalistic settings, such as during everyday activities \cite{Kasneci2017} or using virtual reality environments \cite{Xu2024}.
    
    In this study, we employed the Tobii Pro Glasses 2 system \cite{Onkhar2023} to record binocular eye movements and pupillary responses during the cooking task. The glasses use corneal reflection, paralax compensation to perform  dark-pupil tracking with infrared illumination and a sampling rate of 100 Hz, providing high-precision 2D and 3D gaze position estimates in the scene camera view (25 Hz). Eyer-tracker was calibrated for each participant using a standard 1-point calibration procedure prior to the cooking task. The scene camera records the participant’s egocentric view with a field of view of approximately 82° horizontally and 52° vertically (note that this is narrower than the wearer’s full visual field, which exceeds 160° horizontally and 70° vertically) \cite{TobiiProGlasses2Manual}. Video is captured at a resolution of 1920×1080 pixels, enabling precise mapping of gaze points onto the visual environment.

    Using the event classification framework provided by the Tobii Pro Controller software, we extracted primary gaze events (fixations and saccades) as well as temporal sequences of pupil diameter changes \cite{Onkhar2023}. 
    
    
    \paragraph{Preliminary analysis and future work.\\}

    In concordance with previous study of oculometry during the ATCT task, \cite{10.1145/3486812.3486841}, preliminary analysis suggest that the eye-tracking data quality is sufficient for extracting oculomotor and pupillometric features, although some participants exhibited higher rates of data loss due calibration drift or misplaced glasses (e.g., when superposed to correction's glasses). Moreover, eye-tracking data collected in naturalistic ecological tasks is prone to missing values, due to blinks, gaze loss, or rapid head movements. These are typically non-random and action-dependent, which necessitates careful treatment during feature extraction and downstream analyses \cite{Wolf2023}. Common strategies include interpolation of short gaps, exclusion of low-confidence events, or explicit modelling of missingness patterns. A few illustrative examples of gaze dynamics is reported in \Cref{fig:chapter_4_eye_tracking}.
    
    \begin{figure}[h]
        \centering
        \includegraphics[width=1\linewidth]{chapters/chapter 4/figures/chapter_4_gaze_estimates.png}
        \caption{\textbf{Examples of illustrative gaze dynamics during the cooking task.} (a, b, c) The figure displays gaze dynamics (from blue to red dots) overlaid on the egocentric video frames, showing patterns of fixation (a, $t_0$), saccades between objects (c), gaze anticipation over hand reaching movements (a, c) and attraction of gaze towards the center during head turning (b). Preliminary studies suggested differences in recipe reading patterns (d) and visual exploration strategies (e) between healthy controls and patients with dysexecutive syndromes. These patterns provide insights into the participant's visual attention and cognitive processes during the cooking task.}
        \label{fig:chapter_4_eye_tracking}
    \end{figure}

    Future work will focus on extracting a comprehensive set of gaze, oculomotor, and pupillometric features that capture attentional allocation, visual exploration strategies, and cognitive load during the cooking task. Potential features include \cite{Laborde2025}:

    \begin{itemize}
        \item Fixation metrics: mean and variance of fixation durations, fixation counts, spatial distribution of fixations, and transitions between task-relevant areas of interest (AOIs) such as ingredients, utensils, and stove.
        \item Saccadic metrics: mean and variance of saccade amplitudes, peak velocities, and frequencies, as well as scanpath characteristics (e.g., entropy, fractal dimension).
        \item Pupillometric metrics: baseline pupil diameter, task-evoked pupillary responses (TEPRs) during key recipe steps, and pupil variability measures (e.g., coefficient of variation).
    \end{itemize}
    
    In particular, the temporal alignment of these features with the symbolic representations extracted from egocentric video (\Cref{sec:chapter_5}) will enable to investigate how oculometric and pupillometric dynamics relate to specific action sequences, planning efficiency, and, ultimately, assess executive functioning during the cooking task.


\subsection{From landmarks estimates to tracking, subject re-identification, and application to intervention detection}

    Hands are among the most informative regions for understanding human activity in vision-based behavioral analysis. Fine motor actions are tightly coupled with cognitive control and task planning, making hand movements a key observable in behavioral science \cite{Luxem2023, Friedrich2025}. However, detecting hands in unconstrained settings remains challenging. Hands exhibit substantial variation in shape and scale, are frequently occluded or self-occluded, and often lack the high-contrast visual features (such as the eyes or mouth) that facilitate facial landmark detection. These difficulties are compounded by variable lighting conditions, cluttered backgrounds, and rapid hand motion, all of which demand robust and generalizable models. Recent advances in deep learning–based hand pose estimation have led to significant improvements in hand landmark detection, tracking, gesture recognition, and action anticipation \cite{Friedrich2025, Chen2020}. A growing body of work specifically addresses egocentric (first-person) video content, where hand-centric visual cues are more prominent \cite{ohkawa2023assemblyhandsegocentricactivityunderstanding, Bandini2023}. 

    During the cooking task, participants perform a range of goal-directed actions, involving grasping utensils, breaking chocolate by hand, beating eggs, or setting the oven, offering fine-grained motor patterns that can be quantitatively analyzed from hand landmark tracking. Among the four available cameras, only two provide consistent and informative views of participants’ hands throughout the administration. The rear-facing camera (GoPro $\#1$) rarely captures hands, while the side-view camera (GoPro $\#3$) is excluded as it mostly captures the lateral edge of the hands, which poses significant perspective challenges for landmark detection (e.g., self-occlusion, missing fingers, partial visibility). As a result, we retain only the head-mounted camera and the closest side-facing camera (GoPro2) for hand pose detection.

    We use the open-source MediaPipe Hands model \cite{lugaresi2019mediapipeframeworkbuildingperception} to estimate hand landmarks. For each frame, the pipeline detects up to four hands and returns 21 three-dimensional landmarks per hand, as well as a binary handedness classification with an associated confidence score. % Here we must (for the nature paper) precice the sampling strategy (not use every frame)

    \begin{figure}[!htbp]
        \centering
        \includegraphics[width=1\linewidth]{chapters/chapter 4/figures/chapter_4_hand_landmarks.png}
        \caption{\textbf{Illustrative example of the multi-view hand landmarks estimation.} Hand landmarks estimates are shown from the egocentric recording perspective (top) and from the second fixed camera (bottom). The estimated landmarks are visualized in a 3D coordinate system, with each landmark represented by a colored dot. As indicated by the red rectangle, extreme hand pose or occlusion may lead to missing landmark estimates.}
        \label{fig:chapter_4_hand_landmarks}
    \end{figure}

    One of the key design features of the MediaPipe Hands model to address the challenges highlighted above is its two-stage pipeline. First, a palm detection module, implemented as a lightweight single-shot detector, identifies candidate hand regions within the full-resolution image. This detector outputs oriented bounding boxes that encode both position and coarse rotation of the hand, which are then used to crop and normalize the input region for the second stage. The landmark estimation network is subsequently applied to these cropped regions, predicting 21 three-dimensional hand landmarks per detected hand. The network operates in a regression framework, directly estimating normalized x, y coordinates and depth (z) relative to the wrist joint, enabling precise localization even under complex articulations and occlusions. By separating coarse hand localization from fine-grained landmark regression, the pipeline focuses network capacity on accurate keypoint estimation, reducing the need for extensive geometric data augmentation (e.g., scaling, rotation, translation) while maintaining real-time performance on CPU and mobile devices \cite{lugaresi2019mediapipeframeworkbuildingperception}.

\paragraph{Preliminary analysis and future work.\\}

A qualitative illustration of the resulting hand pose estimation is shown in \Cref{fig:chapter_4_hand_landmarks}. Preliminary analysis led by Alexander Kreibich (who interned at the Borelli Center during the summer of 2024) focused on tracking hand landmarks across contiguous frames to identify manipulation patterns and examiner interventions from the egocentric perspective. While tracking the hands across frames using geometrical constraints and post-hoc morphological filters (opening and closing) stabilize handedness identification, and reduce false positive hands detections, distinguishing between the participant’s and the examiner’s hands remains challenging. As a first step, we propose to leverage the detection of multiple right or left hands to identify candidate examiner intervention instances. Future work will investigate the benefice of combining these candidate regions with audio analysis of communication patterns (presented below), and the contextual information brought by egocentric vision-based symbolic representations, presented in \Cref{sec:chapter_5}.
Future work will explore the integration of these candidate regions with audio-based analyses of communication patterns (presented below), as well as with contextual cues derived from the symbolic representations extracted from egocentric video, discussed in \Cref{sec:chapter_5}.


% Talk about manual evaluation ? 

%  \begin{figure}[!htbp]
%     \centering
%     \includegraphics[width=.6\linewidth]{chapters/chapter 4/figures/chapter_4_intervention_detection.png}
%     \caption{\textbf{Illustrative example of the multi-view hand landmarks estimation.}}
%     \label{fig:chapter_4_intervention_detection}
%  \end{figure}

%  In clinical and cognitive research, reliable hand tracking can help infer motor function, detect tremor, and assess fine motor skills — with applications in neurology, geriatrics, and developmental psychology \cite{Friedrich2025}. Modern algorithms now achieve accuracies comparable to those of gold-standard motion capture systems for many tasks \cite{Friedrich2025}, making them a compelling option for unobtrusive, video-based assessment.


% Detecting hands is a decidedly complex task: our model has to work across a variety of hand sizes with a large scale span ($\approx 20 \times$) and be able to detect occluded and self-occluded hands. Whereas faces have high contrast patterns, e.g., around the eye and mouth region, the lack of such features in hands makes it comparatively difficult to detect them reliably from their visual features alone.

	
% 1. Providing the accurately cropped palm image to the hand
% landmark model drastically reduces the need for data augmentation (e.g. rotations, translation and scale) and allows the network to dedicate most of its capacity towards landmark localization accuracy.

% \cite{Friedrich2025}

% \cite{lugaresi2019mediapipeframeworkbuildingperception} mediapipe 
% \cite{ohkawa2023assemblyhandsegocentricactivityunderstanding}  AssemblyHands: Towards Egocentric Activity Understanding via 3D Hand Pose Estimation The dataset includes synchronized egocentric and exocentric images sampled from the recent Assembly101 dataset, in which participants assemble and disassemble take-apart toys. To obtain high-quality 3D hand pose annotations for the egocentric images, we develop an efficient pipeline, where we use an initial set of manual annotations to train a model to automatically annotate a much larger dataset. Our annotation model uses multi-view feature fusion and an iterative refinement scheme, and achieves an average keypoint error of 4.20 mm, which is 85\% lower than the error of the original annotations in Assembly101. AssemblyHands provides 3.0M annotated images, including 490K egocentric images,

\subsection{Body landmarks extraction}

Locomotion, dynamic balance and control, and postural stability are fundamental components of motor function, closely linked to cognitive processes such as attention, executive function, and spatial awareness \cite{vanderVeen2023, Joseph2020}. For example, recent findings suggest the links between inhibitory control and balance performances \cite{Divandari2024}, and between altered visual attention and balance impairments in patients with RIL \cite{Bargiotas2019}. These motor deficits can significantly impact daily functioning and quality of life.

Body pose estimation from monocular RGB images has seen significant advancements in recent years, driven by deep learning techniques and large-scale annotated datasets \cite{Tang2024AutomaticT}. These methods detect keypoints corresponding to major joints (e.g., head, shoulders, elbows, wrists, hips, knees, ankles) and connect them to form a skeletal representation of the human body. Despite these advances, monocular approaches still face depth ambiguity, self-occlusions, and complex articulations, which can limit the accuracy of 3D pose reconstruction in unconstrained environments \cite{Tang2024AutomaticT}.

In our study, we leveraged the open-source MediaPipe Pose model \cite{lugaresi2019mediapipeframeworkbuildingperception} to extract body landmarks from the rear camera only (Camera \#1 in \Cref{fig:chapter_4_CT_thumbnails}), which provided a consistent and informative viewpoint of participants during the cooking task. The model provides 33 three-dimensional landmarks per detected person, encompassing the head, torso, and limbs, along with a confidence score for each keypoint. Despite relying on a single monocular RGB view, this configuration offers sufficient spatial information to capture the global posture and gross motor dynamics of participants. The resulting skeletal trajectories enable the analysis of static and dynamic balance, and locomotor stability, thereby contributing to the multimodal characterization of participants’ executive and motor functions in an ecologically valid cooking environment.

The MediaPipe pipeline operates in a two-stage manner: a region-of-interest detection module first identifies a coarse bounding box around the person, which is then passed to a regression network that predicts normalized 3D joint locations relative to a root landmark (the hips for the MediaPipe Pose model). By separating coarse localization from fine-grained landmark regression, the model maintains robustness to partial occlusions while providing high temporal resolution suitable for downstream behavioral encoding. The resulting skeletal trajectories can subsequently be used to quantify kinematic features, such as joint angles, movement smoothness, and postural sway, enabling a multimodal assessment of participants’ motor and executive capacities in a naturalistic cooking environment.

\paragraph{Preliminary analysis and future work.\\}

An illustration of body pose estimation is presented in \Cref{fig:chapter_4_body_landmarks}. We observed that the MediaPipe Pose model provides robust and temporally consistent body landmark estimates across a variety of poses and movements. We note, however, that occlusion of participants' arms prevents the analysis of fine motor actions from body landmarks alone, underscoring the importance of integrating hand pose data from the egocentric camera. Nevertheless, the body landmarks capture gross motor patterns (e.g., walking, reaching, or bending to use the oven). We note, however, that body detection may fail intermittently, as can be seen in the figure when the examiner is seated and partially occluded, or during interactions. These may be partially alleviated via landmarks tracking, in a similar spirit than the hand landmarks tracking presented above. Alternatively, recent reviews of pose estimation algorithms \cite{MEDRANOPAREDES2025217} suggest that the MMPose deep learning method could constitute a strong alternative to leverage body movements temporal dynamics \cite{mmmpose}. We discuss strategies to integrate body landmarks with other modalities in the perspectives section (\Cref{sec:chapter_6_multimodal_integration}). In particular, the symbolic representation of participants’ egocentric vision, computed in \Cref{sec:chapter_5}, aim to enrich these analysis through the incorporation of contextual information, such as the progression in the task.

\begin{figure}[!htbp]
    \centering
    \includegraphics[width=1\linewidth]{chapters/chapter 4/figures/chapter_4_body_landmarks.png}
    \caption{\textbf{Illustrative example of the body landmarks estimation.} The figure illustrates the body landmarks estimation for a participant, where the body landmarks are detected and represented in a 3D coordinate system. While body landmarks enable the analysis of gross motor patterns, they may fail to detect fine motor actions due to participant's arms occlusions. Additionally, false negative detections may intermittently occur, for example when the examiner is seated and partially occluded, or during interactions (purple boxes).}
    \label{fig:chapter_4_body_landmarks}
\end{figure}



\subsection{Speech transcription and speaker diarization}
\label{sec:chapter_4_speech_transcription}

Communication patterns and verbal behavior contribute to the assessment of executive functions,  in particular in ecological tasks that engage planning, decision-making, and adaptation to a new environment. During the task, cameras passively record the audio stream while participants are instructed to autonomously follow the recipe, in the presence of the examiner. Patients with DS may exhibit a variety of verbal or communication behaviors, including increased self-directed speech, more frequent requests for help, speech and language impairments, and conversational patterns (e.g., turn-taking and latency of response), which can act as indirect markers of cognitive control, planning deficits, or inhibition failure \cite{MartnezNicols2021,DeLooze2021, Konig2018, Toth2018}.

Recent research has highlighted the relevance of digital markers of both acoustic and linguistic speech features for characterizing neurological impairments: alterations in articulation rate, uncertainty in vowel prediction \cite{Isaev2023Ataxic}, pause frequency and duration, hesitation ratio, pitch variability, jitter, and loudness \cite{MartnezNicols2021, DeLooze2021, Konig2018, Toth2018}.

Until recently, extracting such information from real-world audio recordings required labor-intensive manual transcription. However, advances in large-scale training self-supervised of speech models have significantly improved automatic speech recognition (ASR) and speaker diarization accuracy \cite{baevski2020wav2vec}. In particular, the Whisper model \cite{whisper} and its extended version WhisperX \cite{bain2022whisperx} achieve strong zero-shot transcription performance across multiple languages — including French — and integrate neural speaker diarization by clustering vocal embeddings into speaker-consistent segments. 

\paragraph{Preliminary analysis and future work.\\}

While a detailed implementation with testable demo \cite{rvad_fast} of a recent voice activity detection (VAD) algorithm \cite{TAN20201} was proposed, we opted for using the WhisperX model in this sutdy, which integrates VAD as part of its diarization pipeline. We processed the audio streams recorded from the rear-facing camera in order to extract transcriptions and identify speech segments attributable to either the participant or the examiner, as illustrated in \Cref{fig:chapter_4_speech}. Our qualitative inspection confirms that WhisperX provides robust and temporally accurate transcriptions for downstream speech and communication analysis. 

 \begin{figure}[!htbp]
    \centering
    \includegraphics[width=1\linewidth]{chapters/chapter 4/figures/chapter_4_speech.png}
    \caption{\textbf{Workflow of our preliminary audio analysis.} We applied the state-of-the-art WhisperX speech transcription model to audio recordings captured by the rear-facing camera \cite{bain2022whisperx}. The visualization of a few videos with overlaid speech predictions was used to qualitatively assess the performance of modern speech processing algorithms in the context of Smartflat study.}
    \label{fig:chapter_4_speech}
 \end{figure}

 In particular, future work will include the detection of dialogues to support the automatic detection of examiner intervention,  which has been associated with increased severity of DS \cite{Chevignard2008}, and will be used to complement the hand landmarks–based intervention detection presented above. In addition, extracted speech will be temporally aligned with the symbolic representations extracted from egocentric video (\Cref{sec:chapter_5}) to investigate how verbal behavior relates to specific action sequences, planning efficiency, and executive functioning during the cooking task.



    
    % These variables can be interpreted in light of previous literature:
	% •	Fixation duration and location may reflect attentional allocation and planning efficiency; for example, prolonged fixations or frequent re-fixations on task-relevant AOIs (ingredients, utensils, stove) could indicate increased cognitive effort or reduced executive flexibility \cite{EraslanBoz2024VisualSearch}.
	% •	Saccadic metrics (amplitude, velocity, frequency) provide indices of visual exploration strategies and inhibitory control; altered saccadic patterns have been reported in MCI and in patients with dysexecutive syndromes during tasks requiring monitoring of multiple information streams \cite{Qi2024AbnormalEyeMovements}.
	% •	Pupillometry offers a complementary measure of cognitive load and arousal. Task-evoked pupillary responses (TEPRs) during key moments of the cooking task (e.g., recipe steps requiring working memory or dual-task coordination) may reveal differential engagement of attentional and executive resources between participants with preserved and impaired executive function \cite{Opwonya2024TEPR}.



    % In the context of our cooking task, potential applications of these features include:
    %     1.	Monitoring gaze transitions between task-relevant areas to quantify planning efficiency, sequencing, and task-switching behavior.
    %     2.	Quantifying fixation-saccade patterns to identify altered visual exploration strategies that may reflect executive dysfunction.
    %     3.	Analyzing TEPRs during cognitively demanding recipe steps as an objective marker of cognitive effort and attentional engagement.


%\section{Conclusions}

